%%%%%%%%%%%%%%%%%%%%%%%% Title Page %%%%%%%%%%%%%%%%%%%%%%%%
\documentclass[12pt, a4paper, titlepage]{article}
\usepackage[a4paper,left=2.54cm,right=2.54cm,top=2.54cm,bottom=2.54cm]{geometry}

\usepackage{placeins}
\usepackage{float}
\usepackage{hyperref}
\usepackage{graphicx}
\usepackage{appendix}
\usepackage{setspace}
\usepackage{dcolumn}
\usepackage[printonlyused, withpage]{acronym}
\usepackage{amsmath}
\DeclareMathOperator*{\argmax}{arg\,max} 
\usepackage{natbib}
\makeatletter
\newcommand{\MSonehalfspacing}{%
  \setstretch{1.44}%  default
  \ifcase \@ptsize \relax % 10pt
    \setstretch {1.448}%
  \or % 11pt
    \setstretch {1.399}%
  \or % 12pt
    \setstretch {1.433}%
  \fi
}
\newcommand{\MSdoublespacing}{%
  \setstretch {1.92}%  default
  \ifcase \@ptsize \relax % 10pt
    \setstretch {1.936}%
  \or % 11pt
    \setstretch {1.866}%
  \or % 12pt
    \setstretch {1.902}%
  \fi
}
\makeatother
\MSonehalfspacing


%%%%%%%%%%%%%%DOCUMENT%%%%%%%%%%%%%%
\begin{document}
%%%%%%%%%%%%%%TITLEPAGE%%%%%%%%%%%%%%
\begin{titlepage}
    \begin{center}
    {\LARGE \textbf{Job Title Classification Strategies for the German Labor Market}}
    \\[1cm]
    {\Large \textbf{Masterthesis}}
    \\[1cm]
    {\Large submitted by}
    \\[0.5cm]
    {\LARGE \textbf{Rahkakavee Baskaran}}
    \\[0.5cm]
    {\Large at the}
    \\[0.5cm]
    \includegraphics[width=0.4\textwidth]{logo.jpg}
    \\[1cm]
    {\Large \textbf{Department of Politics and Public Administration}}
    \\[1cm]
    {\Large \textbf{Center for Data and Methods}}
    \\[2cm]
    \begin{minipage}[c]{0.8\textwidth}
    \begin{description}
     \item {\Large \textbf{1.Gutachter:} Prof. Dr. Susumu Shikano}
     \item {\Large \textbf{2.Gutachter:} JunProf Juhi Kulshresthra}
    \end{description}
    \end{minipage}
    \vfill
    {\LARGE \textbf{Konstanz, \today}}
    \end{center}
    \end{titlepage}

%%%%%%%%%%%%%%TableOfContents%%%%%%%%%%%%%%
\tableofcontents
\newpage


%%%%%%%%%%%%%%Abbreviations%%%%%%%%%%%%%%
\section*{Abbreviations}
\begin{acronym}
  \acro{SVM}[SVM]{Support Vector Machine}
  \acro{NB}{NB}(Naive Bayes)
  \acro{LR}{LR}{Logistic Regression}
  \acro{OA}[OA]{overall accuracy}
  \acro{ROC}[ROC]{receiver operating characteristics}
  \acro{TP}[TP]{True positives}
  \acro{TN}[TN]{True negatives}
  \acro{FN}[FN]{False negatives}
  \acro{FP}[FP]{False positives}
  \acro{MAP}[MAP]{maximum a posteriori}
  \acro{KldB}[KldB]{Klassifikation der Berufe 2010}
  \acro{ISCO}[ISCO]{International Standard Clasification of Occupations}
\end{acronym}
\newpage

%%%%%%%%%%%%%%SECTIONS%%%%%%%%%%%%%%
\section{Introduction}
\section{Related work}
\subsection{Textclassification}
Text classification, a highly researched area, is the process of classifying text documents or text segments into a set of predefined classes. During the last decades, researches developed a various number of classifiers. As \cite{Kowsari2019} summarize in their survey of classifiers, we can group the approaches mainly into three groups. The first group contains traditional methods like Naive Bayes (NV), Support Vector Machines (SVM), K-nearest neighbors (KNN),  Logistic Regressions (LR) or Decision Trees (DT) \citep{Vijayan2017, Colas2006, Kowsari2019, Sebastiani2001}. Deep learning methods like Convolutional Neural Networks (CNN) or Recurrent Neural Networks (RNN), which are currently the most advanced algorithms for NLP, form the second group. The last group consists of ensemble learning techniques like Boosting and Bagging.

\subsection{Short Text classification} 
Another potential issue is the length of input documents for classification. Job titles are clearly short text with often not more than 50 characters. Short texts suffer from sparseness, few word co-occurrences, missing shared context, noisiness and ambiguity. Traditional methods, however, are based on word frequency, high word co-occurrence and context, which is why they often fail to achieve high accuracy for short texts \citep{Song2014, WangY2017, WangF2014}. In their overview, \cite{Song2014} present three approaches to solve this. First, since short text data often suffers from unlabeled data in the context of online text data, such as Twitter postings, they suggest using semi-supervised approaches. Second, they recommend to use ensemble learning methods, which focus on the sparseness of the data. Third, \cite{Song2014} propose feature dimensionality reduction and extraction of semantic relationship methods. Based on the latter more recent work on short text classification criticizes the use of the ``Bag of Word'' concept for feature representation as it only reflects the appearance of words in the text. Instead, they represent short texts with semantically similar and conceptual information \citep{Bouaziz2014, WangF2014, Chen2019}.
Another question concerning the representation of short texts is whether to represent them as dense or sparse vectors. In their comparison between tf-idf/counter vectorizer and the dense vectorizer word2vec and doc2vec, \cite{WangY2017} conclude that among the classifiers Naive Bayes, Logisitic Regression and SVM, the sparse vectorizers achieve the highest accuracy. \cite{Chen2019}, conversely, see limitations in sparse representation as it cannot capture the context information. In their work, they integrate sparse and dense representation into a deep neural network with Knowledge powered Attention, which outperform state-of-art deep learning methods, like CNN, for Chinese short texts. 
Concerning the classifiers, there is no consensus approach for short text classification. For traditional approaches \cite{WangY2017}'s results indicate that logistic regression and SVM perform best, while KNN seems to achieve best accuracy in \cite{Khamar2013}'s work. Similar to job title specific work, more recent work prefers deep learning methods, mostly CNN \citep{Chen2019}. 

\subsection{Domain-specific}
Each text classification task presents different challenges. One challenge is that domain-specific problems may arise. There is some work that deals with job classification in the English speaking job market. 
In terms of classifiers, the corresponding work can be categorised into traditional classifiers or deep learning methods. \cite{Zhu2017} for example, use a KNN classifier in combination with document embedding as feature selection strategy. \cite{Javed2015} rely on traditional methods as well, by combining a SVM classifier and a KNN classifier for their job recommendation system. In contrast, the approaches of \cite{Decorte2021}, \cite{WangJ2019} and \cite{Neculoiu2016} are based on Deep Learning methods. From a higher perspective, there is another dividing line between the approaches. As mentioned earlier, job title normalization can be considered as a typical text classification task \citep{WangJ2019, Javed2015, Zhu2017}. \cite{Decorte2021} and \cite{Neculoiu2016}, however, formulate the task as a string representation approach of similar job titles.  

\subsection{Multiclass}
A last challenge of text classification tasks comes with the number of classes. As \cite{Li2004} show in their classification of tissue, multiclass classification is more difficult than binary classification problems. Partly, because most of classification algorithms were designed for binary problems \citep{Aly2005}. Approaches for multiclassification can be grouped into two types. Binary algorithms can handle multiclassification naturally. This is, for example, the case for Regression, DT, SVM, KNN and NV. The second type is the decomposition of the problem into binary classification tasks (for the different subtypes see \cite{Aly2005}). The literature so far does not have a clear answer to solve multiclassification problems. Different approaches, like Boosting \citep{Schapire2000} or CNN \citep{Farooq2017} are applied. It is noticeable, however, that many works use variations of SVM \citep{Guo2015,Tomar2015,Tang2019}. 

\section{Data and taxonomy}

\subsection{KldB 2010 Taxonomy}
The ``\ac{KldB}" is structured hierarchically with 5 levels. On each level there is a different number of classes. On level 1 each class has an id of length one with a number from 0 to 9. Table \ref{tab: T3} shows the 10 classes of level 1 with their class names. On level 2, then, each of the 10 classes are divided into one or more subclasses having a class id of length 2 with the first digit indicating the class of level 1 and the second digit the class of level 2. An overview of the all 5 levels with an example of classes is given in table \ref{tab: T2}. Note that the example in table \ref{tab: T2} does not show on level 2 to level 5 all classes. Thus on level 2 there exists also, e.g. the class id 41 with ``Mathematik-, Biologie- Chemie- und Physikberufe", which in turn is divided into other classes on level 3 etc..  With this procedure this ultimatley leads to class ids of length 5 on level 5. An occupation can be classified on every level in the Taxonomy. Considering the classes of the example in table \ref{tab: T2}, the job title ``Java Developer" could  be classified on level 5 to the class 43412. From this id, it is also derivable that the jobtitle belongs, for example, on level 3 to the class ``Sofwareentwicklung" \citep{Bundesagentur2011a, Bundesagentur2011b, Paulus2013}

The \ac{KldB} contains of two dimension. The first dimension, the so-called "Berufsfachlichkeit" structures jobs according to their similarity in knowledge, activties and jobs. This is reflected in the first 4 levels. Considering the again the example from above and the job title "Fullstack PHP-Entwickler" it is reasonable to classify both on level 1 to "Naturwissenschaft, Geografie and Information", because both of them are related to computer science. It also make sense to classifiy them for example to 4341, because both are about sofware development. On level 5, then, a second dimension is introduced. the "Anforderungsniveau". This dimension gives information on the level of requirement for a job and 4 possible requirements. In table \ref{tab: T4} they are summarized. From the class id of job title ``Java Developer", we can see that the job has been assigned to the second requirement level, since the last digit is a 2 \citep{Bundesagentur2011a,Bundesagentur2011b,Paulus2013}


\begin{table}[]
  \resizebox{\textwidth}{!}{%
  \begin{tabular}{cl}
  \hline
  \multicolumn{1}{l}{\textbf{IDs KldB 2010}} & \textbf{Berufsbereich (Level 1)}                                           \\ \hline
  1                                          & Land-, Forst- und Tierwirtschaft und Gartenbau                             \\
  2                                          & Rohstoffgewinnung, Produktion und Fertigung                                \\
  3                                          & Bau, Architektur, Vermessung und Gebäudetechnik                            \\
  4                                          & Naturwissenschaft, Geografie und Informatik                                \\
  5                                          & Verkehr, Logistik, Schutz und Sicherhe                                     \\
  6                                          & Kaufmännische Dienstleistungen, Warenhandel, Vertrieb, Hotel und Tourismus \\
  7                                          & Unternehmensorganisation, Buchhaltung, Recht und Verwaltung                \\
  8                                          & Gesundheit, Soziales, Lehre und Erziehung                                  \\
  9 & Sprach-, Literatur-, Geistes-, Gesellschafts- und Wirtschaftswissenschaften, Medien, Kunst, Kultur und Gestaltung \\
  0                                          & Militär                                                                    \\ \hline
  \end{tabular}%
  }
  \caption{\label{tab: T3} Overview of classes Level 1 - Berufsbereiche (edited after \citep{Bundesagentur2011b})}
  \end{table}

\begin{table}[hb!]
  \center
  \resizebox{\textwidth}{!}{
  \begin{tabular}{llll}
  \hline
  \textbf{Name}      & \textbf{Level} & \textbf{Number of classes} & \textbf{Example}                               \\ \hline
  Berufsbereiche     & 1              & 10                         & 4: Naturwissenschaft, Geografie und Informatik \\
  Berufshauptgruppen & 2 & 37   & 43: Informatik-, Informations- und Kommunikationstechnologieberufe            \\
  Berufsgruppen      & 3              & 144                        & 434: Sofwareentwicklung                        \\
  Berufsuntergruppen & 4              & 700                        & 4341: Berufe in der Sofwareentwicklung         \\
  Berufsgattungen    & 5 & 1286 & 43412: Berufe in der Sofwareenetwicklung - fachlich ausgerichtete Tätigkeiten \\ \hline
  \end{tabular}%
  }
  \caption{\label{tab: T2} Overview of \ac{KldB} (edited after \citep{Bundesagentur2011b})}
  \end{table}


\begin{table}[]
  \resizebox{\textwidth}{!}{%
  \begin{tabular}{clll}
  \hline
  \multicolumn{1}{l}{\textbf{Level of requirement}} & \textbf{Class ID} & \textbf{Name long} & \textbf{Name short} \\ \hline
  1 & xxxx1 & Helfer- und Anlerntätigkeit        & Helfer     \\
  2 & xxxx2 & fachlich ausgerichtete Tätigkeiten & Fachkraft  \\
  3 & xxxx3 & komplexe Spezialstentätigkeiten    & Spezialist \\
  4 & xxxx4 & hoch komplexe Tätigkeiten          & Experte    \\ \hline
  \end{tabular}%
  }
  \caption{\label{tab: T4} Overview of Level of requirements on Level 5} (edited after \citep{Bundesagentur2011b})
  \end{table}

With the KldB 2010, an useful and information-rich occupational classification was created for Germany that reflects the current trends in the labor market \citep{Paulus2013}. One strength relies in the construction of the \ac{KldB}. Instead of just including expert knowledge into the Taxonomy the developement process is based on systematical consideration of information about occupations, as well as statistical procedures for taxonomy developement. Furthermore, the taxonomy was reviewed qualitatively several times in relation to professions and revised. Considering the expressiveness, the \ac{KldB} has some more benefits. Since the taxonmy is quite recent, it reflects new job classes and market trends very adequately. Further, by including the second dimension, the taxonomy provides a powerful tool to organize job titles into simple requirement classes.  In addition, the taxonomy also distinguishes between managerial, supervisory, and professional employees, which is also valuable information. Finally, the taxonomy also convinces with the possibility to switch to ``\ac{ISCO}" through its IDS and thus to normalize jobs to a global standard \citep{Bundesagentur2011b}

\section{Pipeline}
\section{Preprocessing}
\section{Evaluation metrics}

There exists several metrics for the evaluation of classification approaches in the literature \citep{Fatourechi2008}. The choice of appropriate measurements is a crucial step for obtaining a qualitative comparison in the performance between the baseline algorithms and the new approaches. Often researchers rely on popular metrics like \ac{OA}. However, especially for multiclass and inbalanced dataset tasks it is difficult to rely only on one measure like \ac{OA} In order to select appropriate metrics for comparison in the following the most important metrics will be discussed focussing on multiclass classification and imbalanced data sets. 

Most metrics rely on a confusion matrix. For the multiclass case this confusion matrix is defined as follows \citep{Kautz2017}: 
\begin{table}[hb!]
  \center
  \begin{tabular}{lllll}
  \hline
            & positive examples      &           &             &             \\ \hline
  positive prediction  & $c_{1,1}$ & $c_{1,2}$ & $\dots$     & $c_{1,n}$   \\
            & $c_{2,1}$ & $c_{i,j}$ &             &             \\
            &  $\vdots$         &           & $\ddots$ &   $\vdots$\\
            & $c_{n,1}$ &           & $\dots$     & $c_{n,n}$   \\ \hline
  \end{tabular}
  \caption{\label{tab: T1} Confusion Matrix (edited after \citep[113]{Kautz2017}}
  \end{table}

From the confusion matrix follows that $c_{i, j}$ defines examples which belong to class j and are predicted as class i. Given that $k$ is the current class, \ac{TP} is defined as $tp_{k} = c_{k, k}$, thus examples which are correctly predicted as the current class k. \ac{FN} are defined as those examples which not belonging to the current class k, but are predicted as k. Formally $fn_{k} = \sum_{i=1, i \neq k}^n c_{i, k}$. Next, \ac{TN}, are examples belonging to the current class m, but are not predicted as m. Formally $tn_{k} = \sum_{i=1, i\neq k}^n \sum_{j=1, j \neq k}^n c_{i,j}$. Last, \ac{FP} are defined as examples not belonging to class k, but are predicted as such. Formally this can be expressed as: $fp_{k} =  \sum_{i=1, i \neq k}^n c_{k, i}$ \citep{Kautz2017}

As mentioned the \ac{OA} is one of most common metric for performance evaluation. It represents how well the classifier classifies across all classes correctly. Formally, given that N is number of examples and K the number of all classes, this can be expressed as \citep{Branco2017}: 
\[OA = \frac{1}{K} \sum_{i = 1}^K \frac{tp_{k} + tn_{k}}{N}\]
% \[OA = \frac{TP+ TN}{n}\]
Following the formula an accuracy of 1 means that all examples are correctly classified, while a 0 mean that each example is classified with the wrong class. \citep{Berthold2020}
Although \ac{OA} is a widley used metric it is critized for favouring the majority classes, thus not reflecting minority classes appropriatly in unbalanced datasets \citep{Berthold2020, Fatourechi2008}

Two more popular metrics are precision and recall. Precision respresents how well the classifier detects actual positive examples among the positive predicted examples. Recall, also called sensitivity, in contrast, represents how many examples are labelled as positive among the actual positive examples \citep{Berthold2020}. For the multiclass scenario, two different calculation approaches for each of the metrics are proposed: micro and macro average \citep{Branco2017}. In the macro approach first the metric is calculated for each class k against all other classes. The average of all of them is built. Formally: 

\[precision_{macro} = \frac{1}{K} \sum_{i=1}^k \frac{tp_{i}}{tp_{i} + fp_{i}}\]
\[recall_{macro} = \frac{1}{K} \sum_{i=1}^k \frac{tp_{i}}{tp_{i} + fn_{i}}\]

In contrast the micro approach aggregates the values, which can be formally expressed as follows: 

\[precision_{micro} = \frac{\sum_{i=1}^K tp_i}{\sum_{i=1}^K tp_i + fp_i}\]
\[recall_{micro} = \frac{\sum_{i=1}^K tp_i}{\sum_{i=1}^K tp_i + fn_i}\]

There is a trade-off between precision and recall \citep{Buckland1994}. The F-measure capture both precision and recall by taking the harmonic mean between both. It is calculated as follows \citep{Branco2017,Pan2016}:  

\[F_{micro} = 2 \cdot \frac{precision_{micro} \cdot recall_{micro} }{precision_{micro} + recall_{micro} }\ \]

\[F_{macro} = 2 \cdot \frac{precision_{macro} \cdot recall_{macro} }{precision_{macro} + recall_{macro} }\ \]

Apart from the trade-off between recall and precision, there is also a tradeoff between sensitivity and specificty (1- sensitivity). Using a \ac{ROC}, which plots the specifity against the sensitivity the trade-off can be visualized for different thresholds. The area under the curve then can be used to obtain the performance of the classifier. A large area indicates a better classifier \citep{Berthold2020, Espindola2005}. 

As shown above, there are several metrics for evaluating the performance of a classifier, with the metrics having different focuses. Since the job title classification involves multiclass classification and the descriptive analysis show that the data is clearly unbalanced, at least for some classes in level 5, it is not reasonable to base the evaluation solely on the \ac{OA}. Taking precision, recall and the harmonic mean into account would capture the performance of the minority classes as well. The \ac{ROC} curve does gives, due to its visualization a good impression for the performance, but it is not feasible for high number of classes. Following this argumentation the performance of the classifiers will be evaluated with accuracy, precision, recall, F-measure and Cohen's Kappa. 

\section{Baseline Algorithms}
In order to compare different feature selection methods solid baselines are necessary. As pointed out in the literature review \ac{NB}, \ac{LR} and \ac{SVM} have several advantages for text classifcation tasks. In the following based on a theoretical discussion of each classifier, the exact modeling of the baseline classifiers is justified. 

\subsection{Naive Bayes Classifier}
The \ac{NB}, a family of probabilistic classifiers, uses Bayes' rule in order to determine the most likely class for each document \citep{Schneider2005}. All \ac{NB} classifiers rely on the conditional independence assumptions which means, that ``features are independent of each other, given the category variable'' \citep[48]{Xu2018}. Depending on whether the features are discrete or continous, different distributions, so-called event models are proposed. While from a theoretical perspective for continous features Gaussian distribution is well-suited, for discrete features usually Bernoulli or multinomial distibutions are applied \citep{Xu2018}. Although, popular practical implementations, like the one from sklearn, allow as well fractional counts for multinomial distributions \citep{scikit-learn}. Trying different event models, the multinomial \ac{NB} shows indeed for both Count Vectorizer and for TFIDF the best results, which is why I choose it as event model for the baseline. 

The multinomial \ac{NB} classifies according to the most likely class. Given that a document d has $t = 1, ...., k$ terms and can be assigned to $c = 1,...,j$ classes, the probability of a term in a document given a class is calculated as \citep{Manning2008}:
\[ P(t_k, c_j) = \frac{occurence(t_k, c_j) + 1}{\sum occurence(t, c_j) + |V|} \]

where $|V|$ is the cardinality of the vocabulary. In the denominator 1 is added, so-called Laplace smoothing, in order to avoid zeros, which is the case if the number of terms in a document for one class is zero. \citep{Manning2008}. Further given that $N_c = count(c_j)$ is the number of documents belonging to class $c_j$ and N is number of documents the probability of $c_j$ is defined as $\frac{N_c}{N}$. The probability of a document d belonging to a class $c_j$ can then formulated as follows \citep[258]{Manning2008}:
\[ P(c_j|d) \propto P(c_j) \prod_{i = 1}^k P(t_i|c_j) \]
Then the most likely classes can be determined by \citep{Manning2008}: 
\[\argmax_{c \in C} P(c_j) \prod_{i = 1}^k P(t_i|c_j) \] 

\subsection{Logistic Regression}

\subsection{Support Vector Machines}
However, \ac{SVM} also performed well for text classification. Especially for multiclass tasks, as mentioned in the literature review, often different versions of the algorithm are used and showed good performance \citep{Aiolli2005,Angulo2003,Benabdeslem2006,Guo2015,Mayoraz1999,Tang2019,Tomar2015}. In general \ac{SVM} has several advantages for text classifcation. First, text classifcation usually has a high dimensional input space. \ac{SVM} can handle these large features since they are able to learn independently of the dimensionality of the feature space. In addition \ac{SVM}s are known to perform well for dense and sparse vectors, which is usually the case for text classification \citep{Joachims1998}. Empirical results, for example \citet{Joachims1998} or \cite{Liu2010} confirm the theoretical expectations. It is, therefore, a reasonable option to use a basic version of the \ac{SVM} algorithm as a baseline.

The general idea of a \ac{SVM} is to map ``the input vectors x into a high-dimensional feature space Z through some nonlinear mapping chosen a priori [...], where an optimal separating hyperplane is constructed'' \citep[138]{Vapnik2000}. In \ac{SVM} this optimal hyperplane maximizes the margin, which is simply put the distance from the hyperplane to the closest points, so called Support Vectors, across both classes \citep{Han2012}. Formally, given a training data set with n training vectors $x_i \in R^n, i = 1,....,n$ and the target classes $y_1,...y_i$ with $y_i \in \{-1, 1\}$, the following quadratic programming problem (primal) has to be solved in order to find the optimal hyperplane:
\[\min_{w,b} \frac{1}{2}w^{T}w \] 
\[\text{subject to } y_i(w^T\phi(x_i)+b) \geq 1\]

where $\phi(x_i)$ transforms $x_i$ into a higher dimensional space, $w$ corresponds to the weight and $b$ is the bias \citep{Chang2001,Jordan2006}
The given optimzation function assumes that the data can be separated without errors. This is not always possible, which is why \cite{Cortes1995} introduce a soft margin \ac{SVM}, which allows for missclassfication \citep{Vapnik2000}.
By adding a regularization parameter $C$ with $C > 0$ and the corresponding slack-variable $\xi$ the optimization problem changes to \citep{Chang2001, Han2012}: 
\[\min_{w,b} \frac{1}{2}w^{T}w + C \sum_{i=1}^n \xi_i \] 
\[\text{subject to } y_i(w^T\phi(x_i)+b) \geq 1 - \xi_i, \] 
\[\xi_i \geq, i = 1,...,n\]

Introducing Lagrange multipliers $\alpha_i$ and converting the above optimization problem into a dual problem the optimal $w$ meets \citep{Chang2001, Jordan2006}:
\[w = \sum_{I=1}^n y_i\alpha_i\phi(x_i)\]

with the decision function \citep{Chang2001}:
\[\text{sgn } (w^T\phi(x)+b) = sgn(\sum_{i=1}^n y_i \alpha K(x_i, x) +b)\]

$K(x_i, x)$ corresponds to a Kernel function, which allows to calculate the dot product in the original input space without knowing the exact mapping into the higher space \citep{Han2012, Jordan2006}. 

In order to apply \ac{SVM} to multiclass problems several approaches have been proposed. One stratetgy is to divide the multi-classifcation problem into several binary problems. A common approach here is the one-against-all method. In this method as many \ac{SVM} classifiers are constructed as there are classes k. The k-th classifier assumes that the examples with the k label are positive labels, while all the other examples treated as negative. Another popular approach is the one-against-one method. In this approach $k(k-1)/2$ classifiers are constructed allowing to train in each classifier the data of two classes \citep{Hsu2002}. Besides dividing the multiclass problem into several binary problems, some researches propose approaches to solve the task in one single optimization problem, like \citet{Crammer2001}. \footnote{For a detailed overview of all different methods and the method of \citet{Crammer2001} see \citet{Hsu2002,Crammer2001}}. 

In order to find a strong baseline I checked \ac{SVM}'s with different parameters for the \ac{SVM}, as well as different multiclass approaches. It appears that a SVM using a soft margin with a $C=1$ and a one-vs-rest approach has the best results. I also test different kernels, like  RBF Kernel or linear kernel. The linear kernel, formally  $k(x, x') = x^Tx'$, achieved the best results, which is why I choose it for the baseline. 

\section{Implementation of ...}

\section{Experimental results}
\section{Discussion and Limitations}

\clearpage

\bibliographystyle{apalike}
\bibliography{export}
    
\end{document}